\section{Building Teams}
\label{sec:cses-building-teams}

\problemstatement{Given $n$ players and $m$ friendships, assign each player
to one of two teams so that no two friends are on the same team. Output a
valid assignment, or \texttt{IMPOSSIBLE}.}

\begin{patternbox}
\emph{``Split into two groups so every edge crosses''} is
\textbf{2-colouring}, i.e.\ testing bipartiteness. BFS or DFS colours each
node opposite its neighbour; a conflict (an edge between same-coloured
nodes) means an odd cycle exists and the answer is \texttt{IMPOSSIBLE}.
\end{patternbox}

\subsection*{Two-colour via traversal}

Colour a start node $1$, its neighbours $2$, their neighbours $1$, and so
on, alternating on every edge. If the traversal ever tries to colour a node
the same as an already-coloured neighbour, the graph contains an odd cycle
and cannot be 2-coloured. Repeat over all components, since the graph may
be disconnected.

\begin{ideabox}[Bipartite iff No Odd Cycle]
Alternating colours along edges is consistent exactly when every cycle has
even length. The first same-colour edge encountered is a certificate of an
odd cycle, so the traversal both constructs the teams and detects
impossibility.
\end{ideabox}

\subsection*{Algorithm}

\begin{enumerate}
  \item For each uncoloured node, BFS/DFS colouring neighbours with the
        opposite colour.
  \item If a neighbour already has the same colour, output
        \texttt{IMPOSSIBLE}.
  \item Otherwise output every node's colour ($1$ or $2$).
\end{enumerate}

\subsection*{Dry run}

\dryrun{Friendships $1\!-\!2$, $2\!-\!3$; $n=3$.}

\begin{itemize}
  \item Colour $1{=}1$, $2{=}2$, $3{=}1$. No same-colour edge.
  \item Teams: player $1$ and $3$ on team $1$, player $2$ on team $2$.
\end{itemize}

\subsection*{C++ implementation}

\begin{lstlisting}
#include <bits/stdc++.h>
using namespace std;

int main() {
    int n, m;
    cin >> n >> m;
    vector<vector<int>> adj(n + 1);
    for (int i = 0; i < m; i++) {
        int a, b; cin >> a >> b;
        adj[a].push_back(b);
        adj[b].push_back(a);
    }

    vector<int> color(n + 1, 0);
    for (int s = 1; s <= n; s++) {
        if (color[s]) continue;
        color[s] = 1;
        queue<int> q; q.push(s);
        while (!q.empty()) {
            int u = q.front(); q.pop();
            for (int v : adj[u]) {
                if (!color[v]) { color[v] = 3 - color[u]; q.push(v); }
                else if (color[v] == color[u]) { cout << "IMPOSSIBLE\n"; return 0; }
            }
        }
    }
    for (int i = 1; i <= n; i++) cout << color[i] << " ";
    cout << "\n";
    return 0;
}
\end{lstlisting}

\begin{complexitybox}
\textbf{Time Complexity.} $O(n+m)$ -- one traversal. \textbf{Space
Complexity.} $O(n+m)$.
\end{complexitybox}

\begin{takeawaybox}
Two-team assignment is bipartite 2-colouring: alternate colours along
edges, and a same-colour edge proves an odd cycle. The trick $3-\text{color}$
flips between $1$ and $2$ cleanly. This underlies conflict-graph and
scheduling feasibility checks.
\end{takeawaybox}
